\chapter{Hito 1: Creación del Proyecto}

\section{Objetivo}

El objetivo de este hito es construir la estructura inicial del proyecto y preparar el entorno de desarrollo que utilizaremos durante todo el manual.

Al finalizar este hito el estudiante tendrá:

\begin{itemize}
\item Un repositorio Git inicializado.
\item Una estructura de carpetas profesional.
\item Un entorno virtual de Python.
\item Dependencias básicas instaladas.
\item El primer commit del proyecto.
\end{itemize}

\section{Resultado esperado}

Al finalizar este hito, la estructura del proyecto deberá lucir de la siguiente forma:

\begin{verbatim}
dwh_macrofinanciero_rd/

+-- data/
|   +-- raw/
|   \-- processed/

+-- sql/
+-- etl/
+-- app/
+-- docs/
+-- tests/

+-- README.md
+-- requirements.txt
+-- .gitignore
\end{verbatim}

\section{Paso 1: Crear la carpeta del proyecto}

Abrir PowerShell y navegar al directorio donde se almacenará el proyecto.

Por ejemplo:

\begin{lstlisting}[language=bash]
cd C:\Users\t490\Documents\MacroPol
\end{lstlisting}

Crear la carpeta principal:

\begin{lstlisting}[language=bash]
mkdir dwh_macrofinanciero_rd
cd dwh_macrofinanciero_rd
\end{lstlisting}

Verificar la ubicación actual:

\begin{lstlisting}[language=bash]
pwd
\end{lstlisting}

\section{Paso 2: Abrir el proyecto en VS Code}

Desde la terminal:

\begin{lstlisting}[language=bash]
code .
\end{lstlisting}

VS Code deberá abrir la carpeta recién creada.

\section{Paso 3: Inicializar Git}

Abrir la terminal integrada de VS Code y ejecutar:

\begin{lstlisting}[language=bash]
git init
\end{lstlisting}

Salida esperada:

\begin{verbatim}
Initialized empty Git repository
\end{verbatim}

Verificar el estado:

\begin{lstlisting}[language=bash]
git status
\end{lstlisting}

\section{Paso 4: Crear la estructura de carpetas}

Crear las carpetas principales:

\begin{lstlisting}[language=bash]
mkdir data
mkdir sql
mkdir etl
mkdir app
mkdir docs
mkdir tests
\end{lstlisting}

Crear las subcarpetas para almacenamiento de datos:

\begin{lstlisting}[language=bash]
mkdir data\raw
mkdir data\processed
\end{lstlisting}

\section{Paso 5: Crear los archivos iniciales}

Crear los siguientes archivos:

\begin{itemize}
\item README.md
\item requirements.txt
\item .gitignore
\item .env.example
\end{itemize}

La estructura del proyecto deberá verse así:

\begin{verbatim}
dwh_macrofinanciero_rd/
|
+-- data/
+-- sql/
+-- etl/
+-- app/
+-- docs/
+-- tests/
|
+-- README.md
+-- requirements.txt
+-- .gitignore
+-- .env.example
\end{verbatim}

\section{Paso 6: Configurar el archivo .gitignore}

Crear el archivo:

\begin{verbatim}
.gitignore
\end{verbatim}

con el siguiente contenido:

\begin{lstlisting}
.venv/
.env

**pycache**/

*.pyc
*.pyo

.vscode/

data/processed/

.DS_Store
Thumbs.db
\end{lstlisting}

\section{Paso 7: Crear el README inicial}

Archivo:

\begin{verbatim}
README.md
\end{verbatim}

Contenido sugerido:

\begin{lstlisting}

# Data Warehouse Macrofinanciero RD

Proyecto practico para construir una plataforma
de inteligencia macrofinanciera utilizando:

* PostgreSQL
* Python
* Streamlit
* GitHub

Autor: Francisco Ramirez
\end{lstlisting}

\section{Paso 8: Crear el entorno virtual}

Crear el entorno virtual:

\begin{lstlisting}[language=bash]
python -m venv .venv
\end{lstlisting}

Activarlo:

Windows:

\begin{lstlisting}[language=bash]
.venv\Scripts\activate
\end{lstlisting}

Linux/Mac:

\begin{lstlisting}[language=bash]
source .venv/bin/activate
\end{lstlisting}

La terminal deberá mostrar:

\begin{verbatim}
(.venv)
\end{verbatim}

\section{Paso 9: Instalar dependencias mínimas}

Instalar los paquetes necesarios:

\begin{lstlisting}[language=bash]
pip install pandas
pip install sqlalchemy
pip install psycopg2-binary
pip install python-dotenv
\end{lstlisting}

Guardar las dependencias instaladas:

\begin{lstlisting}[language=bash]
pip freeze > requirements.txt
\end{lstlisting}

\section{Paso 10: Realizar el primer commit}

Agregar todos los archivos al repositorio:

\begin{lstlisting}[language=bash]
git add .
\end{lstlisting}

Verificar:

\begin{lstlisting}[language=bash]
git status
\end{lstlisting}

Crear el primer commit:

\begin{lstlisting}[language=bash]
git commit -m "Initial project structure"
\end{lstlisting}

Consultar el historial:

\begin{lstlisting}[language=bash]
git log --oneline
\end{lstlisting}

Salida esperada:

\begin{verbatim}
3f2a9ab Initial project structure
\end{verbatim}

\section{Checklist de validación}

Antes de continuar al siguiente hito, verificar:

\begin{itemize}
\item[$\square$] VS Code instalado y funcionando.
\item[$\square$] Proyecto creado.
\item[$\square$] Git inicializado.
\item[$\square$] Estructura de carpetas creada.
\item[$\square$] Entorno virtual activo.
\item[$\square$] Dependencias instaladas.
\item[$\square$] Primer commit realizado.
\end{itemize}

\section{Entregable}

Al finalizar este hito el estudiante deberá disponer de:

\begin{itemize}
\item Un proyecto abierto en VS Code.
\item Un repositorio Git funcional.
\item Un entorno virtual configurado.
\item Una estructura base lista para desarrollar el Data Warehouse.
\end{itemize}

\section{Próximo hito}

En el Hito 2 construiremos la infraestructura de base de datos:

\begin{itemize}
\item Instalación de PostgreSQL.
\item Creación de la base de datos.
\item Configuración de variables de entorno.
\item Primera conexión Python $\rightarrow$ PostgreSQL.
\end{itemize}
